\documentclass{article}
\input{../../../lib.sty}

\title{\texttt{fn make\_stable\_group\_by}}
\author{\MichaelShoemateAuthor}
\begin{document}
\maketitle


\contrib

Proves soundness of \rustdoc{transformations/make_stable_lazyframe/group_by/fn}{make\_stable\_group\_by}
in \asOfCommit{mod.rs}{c5a044563}.

\section{Hoare Triple}
\subsection*{Precondition}
\subsubsection*{Compiler-verified}
\begin{itemize}
    \item Generic \texttt{M} implements \rustdoc{transformations/cast_metric/trait}{UnboundedMetric}.
\end{itemize}

\subsubsection*{Caller-verified}
None

\subsection*{Pseudocode}
\label{sec:python-pseudocode}
\lstinputlisting[language=Python,firstline=2,escapechar=|]{./pseudocode/make_stable_group_by.py}

\subsection*{Postcondition}
\begin{theorem}
\validTransformation{\texttt{(input\_domain, input\_metric, plan)}}{\\\texttt{make\_stable\_group\_by}}
\end{theorem}

\begin{proof}[Appropriate Output Domain]
    By line \ref{line:stable-keys}, the execution keys are stable row-by-row expressions.
    The domain analysis on line \ref{line:group-by-domain} excludes floating-point grouping keys.
    Both margin propagation and contribution-bound lookup use only the execution keys.
    The domain analysis requires infallible rewritten aggregates.
    The check also excludes positional inputs and data-dependent quantile
    parameters, ensuring that aggregates depend only on their respective groups.
    The checker admits group-local expressions with data-independent failure behavior.
    Its three certificates are established recursively:
    columns are equivariant vectors; scalar literals and group lengths are invariant scalars.
    Certified elementwise operations preserve the certificate of their input.
    Combining aligned inputs preserves equivariance, with invariant scalars broadcast
    to the vector length; shuffled inputs are never combined this way.
    Filtering and dropping nulls may resize a vector, but do not turn it into a scalar.
    Resizing is forbidden when another argument requires the original alignment;
    a reducer may nevertheless consume a resized vector because its result is scalar.
    A trusted uniform permutation makes an equivariant input invariant in distribution.
    Order-insensitive reducers return invariant scalars, and order-sensitive reducers
    require an invariant scalar or a uniformly shuffled input.
    These rules are conditional on the schema and public-parameter checks for each
    whitelisted Polars operation, which must exclude data-dependent failures.
    The shared preparation inserts permutations for mathematically order-independent
    reducers and certifies exactly the expressions that will execute.
    In particular, the shuffle distribution does not depend on the original
    physical row order.

    By the postcondition of \texttt{DslPlan.simulate\_schema},
    \texttt{series\_domains} follows the expected schema of members in the output domain.
    The output column domains preserve the schema while allowing infallible black-box aggregation.

    Only margins whose grouping expressions are a subset of the execution keys are preserved.
    In the domain analysis, each preserved key becomes \texttt{col(output\_name)}.
    Polars evaluates a grouping expression on the input and stores its value in the output column with that name.
    The translated margin therefore partitions output rows by the same labels as the original margin partitions input rows.
    Each input row contributes to at most one output row,
    so the upper bounds on group lengths and the number of groups remain valid.
    Re-evaluating the original expression after aggregation is not generally equivalent
    and must not inherit the original margin.
    Among these margins, invariants are discarded.

    Therefore the group-by result is in the derived output domain or raises a data-independent runtime exception.
\end{proof}


The stability guarantee doesn't attempt to claim the broadest set of possible bounds on output distances,
rather it only claims a simple bound that might be useful for the user.
This can be extended in the future, but is sufficient for select queries.

\begin{proof}[Stability guarantee]
    We first simplify the problem on line \ref{line:identifier} by only considering datasets that differ by rows, not by identifiers.

    Since ordered input metrics are rejected, row permutations have zero
    distance. The output metric is unchanged, so group ordering is immaterial.
    The only output distance bound derived is for the empty grouping set, namely the entire output table.

    As established by the common domain analysis above, accepted aggregations are invariant
    in distribution under input-row permutations. Their rowwise inputs commute
    with permutations, and order-sensitive reducers require a trusted uniform
    permutation, producing the same distribution regardless of input order.
    For an unchanged group, couple the random permutations at corresponding
    shuffle evaluations so that their shuffled input vectors agree.
    Reuse the same coupling for shared subexpressions, and independent couplings
    for independently evaluated shuffles. This establishes joint invariance of
    the aggregate outputs, rather than merely invariance of each marginal.
    Grouping keys have canonical representatives, so unchanged groups can be
    coupled to have identical output rows.

    We first retrieve optional upper bounds on the number of rows and groups an individual may contribute on lines \ref{line:total-contributions} and \ref{line:total-groups}.
    Both of these bounds directly correspond to the number of rows an individual may influence in the resulting dataset.

    \begin{align}
        &\max_{x \sim x'} d_{\mathrm{FrameDistance<M>}}(\texttt{function}(x), \texttt{function}(x')) \\
        &\le \mathrm{option\_min}(\texttt{contributed\_rows}, \texttt{contributed\_groups}) \cdot 2
    \end{align}

    Since each influenced row accounts for one addition and one removal, the distance is twice the number of influenced rows.
    This is reflected on line \ref{line:double}.
    The resulting bound is returned on line \ref{line:bound}, satisfying the postcondition.
\end{proof}
\end{document}
